\subsection{Ejemplo (Adam)}

% Paleta de curvas de nivel
\colorlet{niv1}{AmarilloAviso!85!black}
\colorlet{niv2}{AmarilloAviso!50!VerdeNeon!75!black}
\colorlet{niv3}{VerdeNeon!65!black}
\colorlet{niv4}{VerdeNeon!40!CelesteBrillante!65!black}
\colorlet{niv5}{CelesteBrillante!65!black}
\colorlet{niv6}{CelesteBrillante!48!black}
\colorlet{niv7}{CelesteBrillante!34!black}

% Comando de referencia adaptado para article
\providecommand{\reflibro}[1]{\par\vspace{0.5em}{\scriptsize\textcolor{GrisTexto!70}{\textit{Ref. libro:} #1}}}

% Curvas de nivel del valle rotado
\newcommand{\adamValle}{%
    \foreach \r/\cl in {0.6/niv1, 1.2/niv2, 1.9/niv3, 2.6/niv4, 3.4/niv5, 4.3/niv6, 5.3/niv7}
        \draw[\cl, rotate around={30:(0,0)}] (0,0) ellipse ({\r} and {\r/3.1623});%
}

% Elipses no rotadas para el ejemplo paso a paso
\newcommand{\momElipses}{%
    \foreach \r/\cl in {0.8/niv1, 1.6/niv2, 2.5/niv3, 3.6/niv4, 5.1/niv5, 6.8/niv6}
        \draw[\cl] (0,0) ellipse ({\r} and {\r/3.1623});%
}

% ---------------------------------------------------------------
% 1. Explicación gráfica: un mismo problema, iteración a iteración
% ---------------------------------------------------------------

\begin{cajaconcepto}{Concepto principal: Adaptive Moment Estimation}
    Mantiene promedios móviles del gradiente y su cuadrado, con una \textbf{corrección de sesgo} inicial:
    \[ v^{(k+1)}=\gamma_v v^{(k)}+(1-\gamma_v)g^{(k)}, \qquad \hat{v}^{(k+1)}=v^{(k+1)}/(1-\gamma_v^k) \]
    \[ s^{(k+1)}=\gamma_s s^{(k)}+(1-\gamma_s)(g^{(k)}\odot g^{(k)}), \qquad \hat{s}^{(k+1)}=s^{(k+1)}/(1-\gamma_s^k) \]
    \[ x^{(k+1)}=x^{(k)}-\alpha\,\hat{v}^{(k+1)}\oslash(\epsilon+\sqrt{\hat{s}^{(k+1)}}) \]
\end{cajaconcepto}

\subsubsection*{Explicación Gráfica - Iteración 1}

\begin{center}
    \begin{tikzpicture}[x=1.5cm, y=1.5cm, font=\small]
        \begin{scope}
            \clip (-3.9,-2.6) rectangle (-0.4,0.2);
            \adamValle
            % GD trace as reference
            \draw[CelesteBrillante, thick, dashed, -stealth] (-3.6,-1.2) -- (-2.43,-1.99);
            % Box movement
            \draw[GrisTexto!60, dashed] (-3.6,-1.2) -- (-3.0,-1.2) -- (-3.0,-1.8);
            % Adam Step
            \draw[NaranjaBase, very thick, -stealth] (-3.6,-1.2) -- (-3.0,-1.8);
        \end{scope}
        \draw[gray!50] (-3.9,-2.6) rectangle (-0.4,0.2);
        \fill[GrisTexto] (-3.6,-1.2) circle (1.8pt) node[above, inner sep=2pt] {$x^{(1)}$};
        \fill[NaranjaBase] (-3.0,-1.8) circle (1.8pt) node[right, inner sep=3pt, text=GrisTexto] {$x^{(2)}$};
        \node[NaranjaBase, anchor=north east] at (-3.2,-1.55) {Adam};
        \node[CelesteBrillante, anchor=south west] at (-2.9,-1.8) {GD};
        \node[anchor=south, text=GrisTexto] at (-3.3,-1.2) {$\alpha$};
        \node[anchor=west, text=GrisTexto] at (-3.0,-1.5) {$\alpha$};
    \end{tikzpicture}
\end{center}

\begin{cajaconcepto}{Inicialización}
    Con $v^{(1)}=0,\ s^{(1)}=0$, el factor de corrección cancela la atenuación inicial:
    \[ \hat{v}^{(2)} = g^{(1)}, \quad \hat{s}^{(2)} = (g^{(1)})^2 \]
\end{cajaconcepto}

\begin{itemize}
    \item Al dividir $\hat{v}$ entre $\sqrt{\hat{s}}$, la magnitud escalar del gradiente se cancela por completo.
    \item El primer paso es aproximadamente $\pm \alpha$ en \textbf{todos} los ejes (movimiento estrictamente diagonal).
    \item Permite ganar inercia rápidamente en las direcciones correctas.
\end{itemize}
\reflibro{Ec.~(5.30)--(5.34).}

\vspace{1em}
\subsubsection*{Explicación Gráfica - Iteración 2}

\begin{center}
    \begin{tikzpicture}[x=1.8cm, y=1.8cm, font=\small]
        \begin{scope}
            \clip (-3.75,-2.55) rectangle (-1.15,-0.75);
            \adamValle
            \draw[NaranjaBase!45, very thick] (-3.6,-1.2) -- (-3.0,-1.8);
            % GD from x2 points perpendicular to the valley
            \draw[GrisTexto!70, thick, dashed, -stealth] (-3.0,-1.8) -- (-3.3,-0.8) node[above, text=GrisTexto!90] {GD};
            % Adam vectors decomposition
            \draw[CelesteBrillante, very thick, -stealth] (-3.0,-1.8) -- (-2.4,-2.4) node[midway, below right, inner sep=1pt] {inercia $\hat{v}$};
            \draw[VerdeNeon, very thick, -stealth] (-2.4,-2.4) -- (-2.4,-1.45) node[midway, right, inner sep=2pt] {escala $1/\sqrt{\hat{s}}$};
            \draw[NaranjaBase, ultra thick, -stealth] (-3.0,-1.8) -- (-2.4,-1.45);
        \end{scope}
        \draw[gray!50] (-3.75,-2.55) rectangle (-1.15,-0.75);
        \fill[GrisTexto!70] (-3.6,-1.2) circle (1.6pt) node[above right, inner sep=2pt] {$x^{(1)}$};
        \fill[GrisTexto] (-3.0,-1.8) circle (1.8pt) node[left, inner sep=3pt] {$x^{(2)}$};
        \fill[NaranjaBase] (-2.4,-1.45) circle (1.8pt) node[right, inner sep=3pt, text=GrisTexto] {$x^{(3)}$};
    \end{tikzpicture}
\end{center}

En esta segunda iteración empezamos a notar el efecto combinado de los dos momentos:
\begin{itemize}
    \item \textcolor{CelesteBrillante}{$\hat{v}$ (Momentum)}: Empuja la búsqueda continuando el rumbo del paso anterior.
    \item \textcolor{VerdeNeon}{$1/\sqrt{\hat{s}}$ (RMSProp)}: Frena agresivamente el movimiento vertical porque el gradiente oscila históricamente allí.
\end{itemize}

\begin{cajaconcepto}{Ajuste dinámico}
    \[ \Delta x_i = -\alpha \frac{\hat{v}_i}{\epsilon + \sqrt{\hat{s}_i}} \]
\end{cajaconcepto}

Como \textbf{resultado}, el paso dobla suavemente dirigiéndose al fondo del valle en lugar de rebotar en sus paredes.

\vspace{1em}
\subsubsection*{Explicación Gráfica - Iteraciones posteriores}

\begin{center}
    \begin{minipage}[c]{0.55\textwidth}
        \centering
        \begin{tikzpicture}[x=1.3cm, y=1.3cm, font=\small]
            \begin{scope}
                \clip (-3.9,-2.6) rectangle (0.6,0.5);
                \adamValle
                % Path of Gradient Descent
                \draw[CelesteBrillante, thick] plot coordinates {(-3.600,-1.200) (-2.406,-2.004) (-2.404,-0.958) (-1.710,-1.289) (-1.619,-0.724) (-1.204,-0.843) (-1.097,-0.530) (-0.842,-0.559) (-0.747,-0.381)};
                % Computed realistic path of Adam
                \draw[NaranjaBase, very thick] plot coordinates {(-3.600,-1.200) (-3.000,-1.800) (-2.400,-1.450) (-1.800,-1.050) (-1.200,-0.700) (-0.700,-0.400) (-0.300,-0.150) (-0.100,-0.050) (0.000,0.000)};
                
                \foreach \p in {(-3.000,-1.800), (-2.400,-1.450), (-1.800,-1.050), (-1.200,-0.700), (-0.700,-0.400), (-0.300,-0.150)}
                    \fill[NaranjaBase] \p circle (1.5pt);
                \foreach \p in {(-2.406,-2.004), (-2.404,-0.958), (-1.710,-1.289), (-1.619,-0.724), (-1.204,-0.843)}
                    \fill[CelesteBrillante] \p circle (1.3pt);
            \end{scope}
            \draw[gray!50] (-3.9,-2.6) rectangle (0.6,0.5);
            \fill[white] (-3.6,-1.2) circle (1.6pt) node[above, inner sep=2pt, text=GrisTexto] {$x^{(1)}$};
            \fill[white] (0,0) circle (1.5pt) node[above left, inner sep=1pt] {$x^*$};
            
            \draw[NaranjaBase, very thick] (-3.8,-2.95) -- (-3.4,-2.95) node[right, text=GrisTexto] {Adam};
            \draw[CelesteBrillante, thick] (-1.3,-2.95) -- (-0.9,-2.95) node[right, text=GrisTexto] {GD};
        \end{tikzpicture}
    \end{minipage}%
    \hfill
    \begin{minipage}[c]{0.4\textwidth}
        \centering
        \begin{tikzpicture}[x=0.6cm, y=2.5cm, font=\footnotesize]
            \draw[gray!60] (0.4,0) -- (7.6,0);
            \foreach \k/\m/\g in {1/0.800/0.632, 2/0.750/0.525, 3/0.600/0.435, 4/0.450/0.361, 5/0.300/0.300, 6/0.150/0.249, 7/0.050/0.207}{
                \fill[NaranjaBase] ({\k-0.32},0) rectangle ({\k},\m);
                \fill[CelesteBrillante] ({\k},0) rectangle ({\k+0.32},\g);
                \node[below] at (\k,0) {$\k$};
            }
            \node[below] at (4,-0.13) {iteración $k$};
            \node[above] at (4,0.9) {Distancia al mínimo};
        \end{tikzpicture}
    \end{minipage}
\end{center}

El algoritmo combina la suavidad estructural de Momentum y el paso regulado por coordenadas de RMSProp:
\begin{itemize}
    \item A diferencia de AdaGrad, el tamaño del paso no decrece prematuramente a cero gracias al decaimiento temporal constante $\gamma_s$.
    \item Evita por completo el zig-zag dramático y destructivo del Gradiente Descendente clásico.
\end{itemize}
\reflibro{§5.8.}

% ---------------------------------------------------------------
% 2. Pseudocódigo
% ---------------------------------------------------------------
\vspace{1em}
\subsubsection*{Pseudocódigo}

\begin{algorithmic}[1]
    \Require $\nabla f$, $x^{(1)}$, $\alpha > 0$, $\gamma_v \in [0,1)$, $\gamma_s \in [0,1)$, tolerancia $\epsilon_{\text{tol}}$, $\epsilon > 0$, $k_{\max}$
    \State $v \gets 0$
    \State $s \gets 0$
    \For{$k = 1, \dots, k_{\max}$}
        \State $g \gets \nabla f(x)$
        \State \textbf{if} $\lVert g \rVert < \epsilon_{\text{tol}}$ \textbf{then break}
        \State $v \gets \gamma_v\, v + (1 - \gamma_v)\, g$ \Comment{(5.30)}
        \State $s \gets \gamma_s\, s + (1 - \gamma_s)\,(g \odot g)$ \Comment{(5.31)}
        \State $\hat{v} \gets v / (1 - \gamma_v^k)$ \Comment{(5.32)}
        \State $\hat{s} \gets s / (1 - \gamma_s^k)$ \Comment{(5.33)}
        \State $x \gets x - \alpha\,\hat{v} \oslash (\epsilon + \sqrt{\hat{s}})$ \Comment{(5.34)}
    \EndFor
    \State \Return $x$
\end{algorithmic}

\vspace{1em}
Los valores típicos sugeridos en la literatura para los hiperparámetros son: $\gamma_v=0.9$, $\gamma_s=0.999$, $\epsilon=10^{-8}$.
\reflibro{Algoritmo~5.9, Ec.~(5.30)--(5.34).}

% ---------------------------------------------------------------
% 3. Ejemplo paso a paso
% ---------------------------------------------------------------
\vspace{1em}
\subsubsection*{Ejemplo Paso a Paso}

\begin{cajaejemplo}{Planteamiento del problema}
    Sea $f(x)=\tfrac12x_1^2+5x_2^2$, donde $\nabla f=(x_1,\ 10x_2)$.
    
    Condiciones iniciales: $x^{(1)}=(-4,\,1)$, $\alpha=0.3$, $\gamma_v=0.9$, $\gamma_s=0.999$.
\end{cajaejemplo}

\noindent
\begin{minipage}[c]{0.55\textwidth}
    \textbf{Inicialización:} $v=(0,\ 0), \quad s=(0,\ 0)$ \\[0.5em]
    \textbf{Iteración $k=1$} \\
    $g^{(1)}=(-4,\ 10)$ \\
    $v^{(2)}=(1-\gamma_v)g^{(1)}=(-0.4,\ 1.0)$ \\
    $s^{(2)}=(1-\gamma_s)(g^{(1)})^2=(0.016,\ 0.1)$ \\
    $\hat{v}^{(2)}=v^{(2)}/(1-0.9)=(-4,\ 10)$ \\
    $\hat{s}^{(2)}=s^{(2)}/(1-0.999)=(16,\ 100)$ \\
    $x^{(2)}=x^{(1)} - 0.3 \frac{(-4,\ 10)}{\sqrt{(16,\ 100)}} = (-3.70,\ 0.70)$ \\
    $\mathbf{f(x^{(2)})=9.295}$
\end{minipage}%
\hfill
\begin{minipage}[c]{0.4\textwidth}
    \centering
    \begin{tikzpicture}[x=1.2cm, y=1.2cm, font=\small]
        \begin{scope}
            \clip (-4.4,-1.3) rectangle (0.6,1.35);
            \momElipses
            \draw[NaranjaBase, very thick, -stealth] (-4,1) -- (-3.70,0.70);
        \end{scope}
        \draw[gray!50] (-4.4,-1.3) rectangle (0.6,1.35);
        \fill[GrisTexto] (-4,1) circle (1.8pt) node[right, inner sep=3pt] {$x^{(1)}$};
        \fill[GrisTexto] (-3.70,0.70) circle (1.8pt) node[right, inner sep=3pt] {$x^{(2)}$};
        \fill[white] (0,0) circle (1.6pt) node[above, inner sep=2pt] {$x^*$};
    \end{tikzpicture}
\end{minipage}
\vspace{1em}

\noindent
\begin{minipage}[c]{0.55\textwidth}
    \textbf{Iteración $k=2$} \\
    $g^{(2)}=(-3.70,\ 7.00)$ \\
    $v^{(3)}=(-0.73,\ 1.60), \quad \hat{v}^{(3)}\approx(-3.84,\ 8.42)$ \\
    $s^{(3)}\approx(0.03,\ 0.15), \quad \hat{s}^{(3)}\approx(14.85,\ 74.96)$ \\
    $x^{(3)}=(-3.70,\ 0.70) - 0.3 \frac{(-3.84,\ 8.42)}{\sqrt{(14.85,\ 74.96)}}$ \\
    $x^{(3)}=(-3.40,\ 0.41)$ \\
    $\mathbf{f(x^{(3)})=6.612}$
\end{minipage}%
\hfill
\begin{minipage}[c]{0.4\textwidth}
    \centering
    \begin{tikzpicture}[x=1.6cm, y=1.6cm, font=\small]
        \begin{scope}
            \clip (-4.25,-1.75) rectangle (-1.9,1.15);
            \momElipses
            \draw[GrisTexto!60, thick, -stealth] (-4,1) -- (-3.70,0.70);
            \draw[NaranjaBase, very thick, -stealth] (-3.70,0.70) -- (-3.40,0.41);
        \end{scope}
        \draw[gray!50] (-4.25,-1.75) rectangle (-1.9,1.15);
        \fill[GrisTexto] (-4,1) circle (1.6pt) node[right, inner sep=3pt] {$x^{(1)}$};
        \fill[GrisTexto] (-3.70,0.70) circle (1.6pt) node[left, inner sep=2pt] {$x^{(2)}$};
        \fill[VerdeNeon] (-3.40,0.41) circle (1.8pt) node[right, inner sep=2pt, text=GrisTexto] {$x^{(3)}$};
    \end{tikzpicture}
\end{minipage}
\vspace{1em}

\noindent
\begin{minipage}[c]{0.55\textwidth}
    \textbf{Iteración $k=3$} \\
    $g^{(3)}=(-3.40,\ 4.07)$ \\
    $x^{(4)}=(-3.10,\ 0.13)$ \\
    $\mathbf{f(x^{(4)})=4.902}$
\end{minipage}%
\hfill
\begin{minipage}[c]{0.4\textwidth}
    \centering
    \begin{tikzpicture}[x=1.1cm, y=1.1cm, font=\small]
        \begin{scope}
            \clip (-4.4,-1.15) rectangle (0.5,1.3);
            \momElipses
            \draw[NaranjaBase, very thick] plot coordinates {(-4.00,1.00) (-3.70,0.70) (-3.40,0.41) (-3.10,0.13) (-2.81,-0.11) (-2.52,-0.30) (-2.23,-0.42)};
            \foreach \p in {(-3.70,0.70), (-3.40,0.41), (-3.10,0.13), (-2.81,-0.11), (-2.52,-0.30), (-2.23,-0.42)}
                \fill[NaranjaBase] \p circle (1.4pt);
        \end{scope}
        \draw[gray!50] (-4.4,-1.15) rectangle (0.5,1.3);
        \fill[white] (-4,1) circle (1.6pt) node[right, inner sep=3pt, text=GrisTexto] {$x^{(1)}$};
        \fill[white] (0,0) circle (1.4pt);
        \draw[NaranjaBase, very thick] (-4.3,-1.45) -- (-3.9,-1.45) node[right, text=GrisTexto] {Adam};
    \end{tikzpicture}
    
    {\footnotesize Primeras 7 iteraciones de Adam.\par}
\end{minipage}

\vspace{1.5em}
\begin{cajaejemplo}{Resultado final ($k_{\max}=60$)}
    Tras varias iteraciones, el algoritmo estabiliza su descenso: 
    \quad $x^{(61)} = (-0.0158,\ 0.0135)$ \quad $\mathbf{f(x^{(61)}) = 0.0010}$
\end{cajaejemplo}

\vspace{1.5em}
\begin{center}
    \begin{tabular}{l c c c c c c c}
        \toprule
        $k$ & 0 & 1 & 2 & 3 & 4 & 5 & 6 \\
        \midrule
        $f(x^{(k+1)})$ & 13.00 & 9.295 & 6.612 & 4.901 & 4.002 & 3.614 & 3.380 \\
        \bottomrule
    \end{tabular}
\end{center}

% ---------------------------------------------------------------
% 4. Código
% ---------------------------------------------------------------
\vspace{1em}
\subsubsection*{Implementación en Python}

\begin{lstlisting}[language=Python]
import numpy as np

def adam(grad, x, alpha=0.001, gamma_v=0.9, gamma_s=0.999, 
         eps=1e-8, k_max=1000, eps_tol=1e-6):
    v = np.zeros_like(x)                  # inicializacion del primer momento
    s = np.zeros_like(x)                  # inicializacion del segundo momento
    
    for k in range(1, k_max + 1):
        g = grad(x)
        if np.linalg.norm(g) < eps_tol:   # criterio de parada
            break
            
        v = gamma_v * v + (1 - gamma_v) * g
        s = gamma_s * s + (1 - gamma_s) * (g**2)
        
        v_hat = v / (1 - gamma_v**k)      # correccion de sesgo para v
        s_hat = s / (1 - gamma_s**k)      # correccion de sesgo para s
        
        x = x - alpha * v_hat / (np.sqrt(s_hat) + eps) # (5.34)
        
    return x, k

# Ejemplo: f(x) = 0.5*x1^2 + 5*x2^2   ->   grad f = (x1, 10*x2)
grad = lambda x: np.array([x[0], 10*x[1]])
x0 = np.array([-4.0, 1.0])
adam(grad, x0, alpha=0.1, k_max=100)      # convergencia del ejemplo
\end{lstlisting}

% ---------------------------------------------------------------
% 5. Análisis de convergencia
% ---------------------------------------------------------------
\vspace{1em}
\subsubsection*{Análisis de Convergencia}

\begin{cajaalerta}{Sensibilidad a hiperparámetros}
    Adam adapta pasos por coordenada, pero la \textbf{tasa base $\alpha$} sigue dictando la estabilidad global. En la práctica empírica, ajustar $\alpha$ suele bastar mientras $\gamma_v, \gamma_s$ se mantienen en sus valores por defecto.
\end{cajaalerta}

\vspace{1em}
\begin{center}
    \begin{minipage}[c]{0.55\textwidth}
        \centering
        \begin{tikzpicture}[x=0.08cm, y=0.3cm, font=\scriptsize]
            % Ejes y fondo
            \draw[gray!60] (0,-12) -- (60,-12);
            \draw[gray!60] (0,-12) -- (0,2);
            \foreach \yy/\lb in {0/{10^{0}}, -4/{10^{-4}}, -8/{10^{-8}}, -12/{10^{-12}}}{
                \draw[gray!60] (-0.8,\yy) -- (0,\yy);
                \draw[gray!25] (0,\yy) -- (60,\yy);
                \node[left] at (-0.8,\yy) {$\lb$};
            }
            \foreach \xx in {0,10,...,60}
                \draw[gray!60] (\xx,-12) -- (\xx,-12.4) node[below] {$\xx$};
            \node[below] at (30,-13.3) {iteración $k$};
            \node[rotate=90] at (-12,-5) {$f(x^{(k)})-f^*$};
            
            % Trazos simulados para Adam con diferentes alpha
            \draw[CelesteBrillante, thick] plot[smooth] coordinates {(0,1.11) (10,0.9) (20,0.6) (30,0.2) (40,-0.2) (50,-0.6) (60,-1.0)};
            \draw[RojoAlerta, thick] plot[smooth] coordinates {(0,1.11) (5,0.0) (10,0.5) (15,-0.2) (20,0.8) (25,-0.5) (30,1.0) (35,-0.1) (40,1.5) (45,0.2) (50,1.8) (55,0.5) (60,2.0)};
            \draw[NaranjaBase, very thick] plot[smooth] coordinates {(0,1.11) (5,0.5) (10,0.1) (15,0.15) (20,-0.32) (25,-0.10) (30,-0.48) (40,-1.44) (50,-1.38) (60,-3.00)};
            
            % Leyenda
            \node[anchor=west, text=CelesteBrillante] at (35,0.6) {$\alpha=0.05$ (lento)};
            \node[anchor=west, text=RojoAlerta] at (42,2.5) {$\alpha=1.0$};
            \node[anchor=west, text=NaranjaBase] at (33,-2.2) {$\alpha=0.3$};
        \end{tikzpicture}
        
        \vspace{0.5em}
        {\footnotesize Comportamiento de Adam variando la tasa de aprendizaje $\alpha$.\par}
    \end{minipage}%
    \hfill
    \begin{minipage}[c]{0.4\textwidth}
        \centering
        \begin{tabular}{l c}
            \toprule
            Tasa $\alpha$ & Comportamiento \\
            \midrule
            Demasiado baja & Lento (subóptimo) \\
            Bien ajustada & \textbf{Rápido y estable} \\
            Demasiado alta & Diverge / Oscila \\
            \bottomrule
        \end{tabular}
    \end{minipage}
\end{center}

\vspace{1.5em}

\noindent
\begin{minipage}[t]{0.48\textwidth}
    \centering
    \textbf{El problema del tamaño de paso efectivo}
    \vspace{1em}
    
    \begin{tikzpicture}[x=1.5cm, y=1.7cm, font=\small]
        % Ejes
        \draw[-stealth, GrisTexto] (0,0) -- (3.5,0) node[right] {$k$};
        \draw[-stealth, GrisTexto] (0,0) -- (0,2.2) node[above, align=center] {Paso efectivo\\$\alpha / \sqrt{\hat{s}^{(k)}}$};
        
        % Curva Adam
        \draw[NaranjaBase, thick] plot[smooth] coordinates {(0.2,1.8) (0.8,0.7) (1.5,1.2) (2.2,0.4) (2.8,0.9) (3.2,0.3)};
        % Curva AMSGrad
        \draw[VerdeNeon, very thick, dashed] plot[smooth] coordinates {(0.2,1.8) (0.8,0.7) (1.5,0.7) (2.2,0.4) (2.8,0.4) (3.2,0.3)};
        
        % Anotaciones
        \node[text=NaranjaBase, anchor=west] at (1.5,1.4) {Adam};
        \node[text=VerdeNeon, anchor=west] at (0.9,0.25) {AMSGrad};
        
        % Puntos destacados
        \fill[NaranjaBase] (1.5,1.2) circle (1.5pt);
        \fill[NaranjaBase] (2.8,0.9) circle (1.5pt);
        \draw[gray!60, dotted] (1.5,0.7) -- (1.5,1.2);
        \draw[gray!60, dotted] (2.8,0.4) -- (2.8,0.9);
    \end{tikzpicture}
    
    \vspace{0.5em}
    {\footnotesize El denominador $\sqrt{\hat{s}}$ fluctúa, permitiendo picos en el tamaño de paso.\par}
\end{minipage}%
\hfill
\begin{minipage}[t]{0.48\textwidth}
    \begin{itemize}
        \item \textbf{Garantías teóricas:} Adam posee un límite de arrepentimiento de $\mathcal{O}(\sqrt{T})$ en problemas convexos.
        \item \textbf{Fallas en convexidad:} Se ha demostrado que Adam \textbf{puede no converger} en problemas simples donde un gradiente grande y esporádico es olvidado rápidamente por el decaimiento $\gamma_s$.
        \item \textbf{Solución (AMSGrad):} Retiene el máximo histórico del segundo momento para forzar un paso monótono decreciente:
        \[ \hat{s}^{(k)} = \max(\hat{s}^{(k-1)}, s^{(k)}) \]
    \end{itemize}
    \vspace{1em}
    \reflibro{Reddi et al. (2018) "On the Convergence of Adam and Beyond".}
\end{minipage}